\chapter{CONSIDERAÇÕES FINAIS}
Este trabalho apresentou o projeto, a construção e a avaliação de uma haste de um grau de liberdade estabilizada por empuxo, empregada como estudo de caso para a aplicação do FreeRTOS a um sistema de controle crítico. As tarefas foram divididas segundo o período e a criticidade: a cadeia sensor, filtro, PID e motor roda em uma única tarefa, acordada a cada 5~ms por um semáforo liberado pelo TIM2, e comandos e telemetria ficam em tarefas de menor prioridade, que não conseguem atrasá-la. Medido no microcontrolador, o atraso entre a interrupção e o PWM foi de no máximo 580~$\mu$s, 11,6\% do período, sem nenhum ciclo perdido. Os módulos, a configuração do STM32 e as versões das ferramentas estão documentados para permitir a replicação.

Parte significativa do aprendizado decorreu dos problemas enfrentados na bancada. O aquecimento do motor foi atribuído à corrente pulsada do PWM sobre um motor de indutância muito baixa e resolvido com um indutor em série. Com a hélice em rotação, o acelerômetro deixou de ser confiável, e o ângulo em voo passou a ser obtido apenas pelo giroscópio; o viés deste, contudo, provocou uma deriva de cerca de $4^\circ$/min e, na ausência de corte por ângulo, a haste chegou a completar uma volta. A solução foi o corte por ângulo e uma correção muito lenta pelo acelerômetro em voo, que reduziu a diferença entre o ângulo do firmware e o do acelerômetro de 13,8$^\circ$ para 1 a 2$^\circ$. O ajuste do PID mostrou que, com atraso de empuxo de dezenas de milissegundos, ganhos altos oscilam e que a integração condicional evita o \textit{windup}; a simulação com a planta identificada quantificou esse efeito.

O refino na bancada, com cinco conjuntos de ganhos, 1\,200 amostras em regime e degraus de \textit{setpoint} em cada um, mostrou que em regime todos são equivalentes (erro quadrático médio de $0{,}29^\circ$ a $0{,}32^\circ$) e que nos degraus $K_i=0{,}08$ reduz o erro acumulado em 28\% em relação aos ganhos de projeto; esse foi o estado final adotado e regravado. A bancada, entretanto, não leva a saída à saturação e, por isso, não exercita o anti-\textit{windup}, cujo efeito foi evidenciado apenas na simulação com ganhos elevados.

Como limitações do trabalho, destacam-se: a medição de tempo foi feita com a telemetria ativa e não prova o pior caso; as curvas de distúrbio e o ensaio do acelerômetro com o motor desligado não foram registrados; o ângulo em voo depende de uma correção lenta pelo acelerômetro, com erro fixo igual à polarização dele, ainda não conferida com um nível externo; não há \textit{watchdog} independente, de modo que uma falha do processador deixa o PWM no último valor. Como trabalhos futuros, sugerem-se: repetir a medição de tempo sem a telemetria, adicionar um \textit{watchdog} e a verificação de estouro de pilha, conferir o zero físico com um nível externo, ler o sensor por DMA e disparar o ciclo pelo pino de dado pronto do MPU6050.
